% Distilled blueprint for Evans, Chapter 6 (Second-Order Elliptic Equations).
% Formal declarations, metadata, and citation identifiers are preserved; exposition and exercises are omitted.

Throughout, $U\subset\R^n$ is a bounded open set and $a^{ij}=a^{ji}$ unless a declaration states otherwise.

\section{Definitions}
\label{sec:elliptic-equations-definitions}

\subsection{Elliptic equations}

\begin{definition}[Divergence and nondivergence form]
\label{def:divergence-nondivergence-form}
\dcref{ch6:6.1-def-1}
\group{Elliptic PDE}
Let
$$(2) \qquad Lu = -\sum_{i,j=1}^n \left(a^{ij}(x)u_{x_i}\right)_{x_j} + \sum_{i=1}^n b^i(x)u_{x_i} + c(x)u$$
or else
$$(3) \qquad Lu = -\sum_{i,j=1}^n a^{ij}(x)u_{x_ix_j} + \sum_{i=1}^n b^i(x)u_{x_i} + c(x)u,$$
for given coefficient functions $a^{ij}, b^i, c$ $(i,j=1,\ldots,n)$.

The PDE is in divergence form for (2) and nondivergence form for (3); $u=0$ on $\partial U$ is the Dirichlet condition.
\end{definition}

\begin{remark}[Converting between divergence and nondivergence form]
\label{rem:elliptic-operator-form-conversion}
\uses{def:divergence-nondivergence-form}
\dcref{ch6:6.1-rem-1}
\group{Elliptic PDE}

If $a^{ij}\in C^1(\bar U)$, the divergence operator (2) equals
$$(2') \qquad Lu = -\sum_{i,j=1}^n a^{ij}(x)u_{x_ix_j} + \sum_{i=1}^n \tilde b^i(x)u_{x_i} + c(x)u$$
with $\tilde b^i := b^i - \sum_{j=1}^n a^{ij}_{x_j}$. Thus the two forms are equivalent when the leading coefficients are $C^1$; divergence form is adapted to energy identities and nondivergence form to pointwise maximum principles.
\end{remark}

\begin{definition}[Uniform ellipticity]
\label{def:uniform-ellipticity}
\uses{def:divergence-nondivergence-form}
\dcref{ch6:6.1-def-2}
\group{Elliptic PDE}

We say the partial differential operator $L$ is (uniformly) elliptic if there exists a constant $\theta > 0$ such that
$$(4) \qquad \sum_{i,j=1}^n a^{ij}(x)\xi_i\xi_j \geq \theta|\xi|^2$$
for a.e. $x \in U$ and all $\xi \in \R^n$.
\end{definition}

\begin{remark}[Physical interpretation]
\label{rem:elliptic-pde-physical-interpretation}
\uses{def:uniform-ellipticity}
\dcref{ch6:6.1-rem-2}
\group{Elliptic PDE}

The principal part $\mathbf A:D^2u=\sum_{i,j}a^{ij}u_{ij}$ models anisotropic diffusion. With flux $\mathbf F=-\mathbf A Du$, ellipticity implies
$$\mathbf F \cdot Du \leq 0;$$
while $\mathbf b\cdot Du$ is transport and $cu$ is a zeroth-order source term.
\end{remark}

\subsection{Weak solutions}

\begin{definition}[Bilinear form and weak solution for $f \in L^2(U)$]
\label{def:bilinear-form-weak-solution-l2}
\uses{def:divergence-nondivergence-form}
\dcref{ch6:6.1-def-3}
\group{Weak Solutions}

(i) The bilinear form $B[\ ,\ ]$ associated with the divergence form elliptic operator $L$ defined by (2) is
$$(8) \qquad B[u,v] := \int_U \sum_{i,j=1}^n a^{ij}u_{x_i}v_{x_j} + \sum_{i=1}^n b^iu_{x_i}v + cuv \, dx$$
for $u,v \in H_0^1(U)$.

(ii) We say that $u \in H_0^1(U)$ is a \textit{weak solution} of the boundary-value problem (1) if
$$(9) \qquad B[u,v] = (f,v)$$
for all $v \in H_0^1(U)$, where $(\ ,\ )$ denotes the inner product in $L^2(U)$.
\end{definition}

\begin{remark}[Variational formulation]
\label{rem:variational-formulation}
\uses{def:bilinear-form-weak-solution-l2}
\dcref{ch6:6.1-rem-3}
\group{Weak Solutions}

Identity (9) is the variational formulation of (1).
\end{remark}

\begin{definition}[Weak solution of the elliptic boundary-value problem]
\label{def:weak-solution-elliptic-bvp}
\uses{def:bilinear-form-weak-solution-l2, thm:characterization-h-minus-1}
\dcref{ch6:6.1-def-4}
\group{Weak Solutions}

We say $u \in H_0^1(U)$ is a weak solution of problem (10) provided
$$B[u,v] = \langle f,v \rangle$$
for all $v \in H_0^1(U)$, where $\langle f,v \rangle = \int_U f^0 v + \sum_{i=1}^n f^i v_{x_i} \, dx$ and $\langle \cdot,\cdot \rangle$ is the pairing of $H^{-1}(U)$ and $H_0^1(U)$.
\end{definition}

\begin{remark}[Nonzero boundary conditions]
\label{rem:nonzero-boundary-weak-solution}
\uses{def:weak-solution-elliptic-bvp, def:trace-operator}
\dcref{ch6:6.1-rem-4}
\group{Weak Solutions}

Assume $\partial U$ is $C^1$ and $u\in H^1(U)$ solves
$$\begin{cases} Lu = f & \text{in } U \\ u = g & \text{on } \partial U. \end{cases}$$
If $g$ is the trace of $w\in H^1(U)$, then $\tilde u:=u-w\in H_0^1(U)$ solves
$$\begin{cases} L\tilde u = \tilde f & \text{in } U \\ \tilde u = 0 & \text{on } \partial U, \end{cases}$$
where $\tilde f := f - Lw \in H^{-1}(U)$.
\end{remark}

\section{Existence of Weak Solutions}
\label{sec:existence-of-weak-solutions}

\subsection{Lax--Milgram Theorem}

\begin{theorem}[Lax--Milgram Theorem]
\label{thm:lax-milgram-theorem}
\dcref{ch6:6.2-thm-1}
\group{Weak Solutions}
Assume that
$$B : H \times H \to \R$$
is a bilinear mapping, for which there exist constants $\alpha,\beta > 0$ such that
$$(i) \qquad |B[u,v]| \leq \alpha\|u\| \, \|v\| \quad (u,v \in H)$$
and
$$(ii) \qquad \beta\|u\|^2 \leq B[u,u] \quad (u \in H).$$
Finally, let $f : H \to \R$ be a bounded linear functional on $H$.

Then there exists a unique element $u \in H$ such that
$$(1) \qquad B[u,v] = \langle f,v \rangle$$
for all $v \in H$.
\end{theorem}

\begin{proof}
\sketch Riesz representation defines a bounded operator $A$ by $B[u,v]=(Au,v)$. Coercivity gives $\beta\|u\|\le\|Au\|$, hence injectivity and closed range; the same inequality excludes a nonzero vector orthogonal to the range, so $A$ is surjective. Riesz representation of $f$ and coercivity then give existence and uniqueness.
\end{proof}

\begin{remark}[The symmetric case]
\label{rem:lax-milgram-symmetric-case}
\uses{thm:lax-milgram-theorem}
\dcref{ch6:6.2-rem-1}
\group{Weak Solutions}

If $B$ is symmetric,
$$B[u,v] = B[v,u] \quad (u,v \in H),$$
then $B$ itself is an inner product equivalent to the Hilbert norm, so Riesz representation gives the result directly. Lax--Milgram does not require symmetry.
\end{remark}

\subsection{Energy estimates}

\begin{theorem}[Energy estimates]
\label{thm:energy-estimates}
\uses{def:bilinear-form-weak-solution-l2, def:uniform-ellipticity, thm:poincare-w01p-subcritical}
\dcref{ch6:6.2-thm-2}
\group{Weak Solutions}

There exist constants $\alpha,\beta > 0$ and $\gamma \geq 0$ such that
$$(i) \qquad |B[u,v]| \leq \alpha\|u\|_{H_0^1(U)}\|v\|_{H_0^1(U)}$$
and
$$(ii) \qquad \beta\|u\|_{H_0^1(U)}^2 \leq B[u,u] + \gamma\|u\|_{L^2(U)}^2$$
for all $u,v \in H_0^1(U)$.
\end{theorem}

\begin{proof}
\sketch Boundedness of the coefficients and Cauchy--Schwarz prove (i). Ellipticity controls $\|Du\|_{L^2}^2$; absorbing the first-order term with Cauchy inequality and using Poincare inequality yields (ii).
\end{proof}

\begin{theorem}[First Existence Theorem for weak solutions]
\label{thm:first-existence-weak-solutions}
\uses{thm:lax-milgram-theorem, thm:energy-estimates, def:bilinear-form-weak-solution-l2}
\dcref{ch6:6.2-thm-3}
\group{Weak Solutions}

There is a number $\gamma \geq 0$ such that for each
$$(7) \qquad \mu \geq \gamma$$
and each function
$$f \in L^2(U),$$
there exists a unique weak solution $u \in H_0^1(U)$ of the boundary-value problem
$$(8) \qquad \begin{cases} Lu + \mu u = f & \text{in } U \\ u = 0 & \text{on } \partial U \end{cases}$$
\end{theorem}

\begin{proof}
\sketch For $\mu\ge\gamma$, the shifted form $B_\mu[u,v]=B[u,v]+\mu(u,v)$ is bounded and coercive. Apply Lax--Milgram to the functional $v\mapsto(f,v)$.
\end{proof}

\begin{remark}[General right-hand side and the isomorphism $L_\mu$]
\label{rem:general-rhs-weak-solution-isomorphism}
\uses{thm:first-existence-weak-solutions, def:weak-solution-elliptic-bvp}
\dcref{ch6:6.2-rem-2}
\group{Weak Solutions}

For
$$f^i \in L^2(U) \quad (i=0,\ldots,n),$$
there is a unique weak solution $u$ of
$$(9) \qquad \begin{cases} Lu + \mu u = f^0 - \sum_{i=1}^n f^i_{x_i} & \text{in } U \\ u = 0 & \text{on } \partial U. \end{cases}$$
because $\langle f,v\rangle = \int_U f^0v + \sum_{i=1}^n f^iv_{x_i}\,dx$ is bounded on $H_0^1(U)$.

In particular, we deduce that the mapping
$$L_\mu = L + \mu I : H_0^1(U) \to H^{-1}(U) \quad (\mu \geq \gamma)$$
is an isomorphism.
\end{remark}

\begin{example}[Nonnegative zeroth-order coefficient]
\label{ex:energy-estimates-gamma-zero-examples}
\uses{thm:energy-estimates, thm:poincare-w01p-subcritical}
\dcref{ch6:6.2-ex-1}
\group{Weak Solutions}

For $Lu=-\Delta u$, the energy estimate holds with $\gamma=0$ by Poincare inequality. The same is true for $Lu=-\sum_{i,j}(a^{ij}u_{x_i})_{x_j}+cu$ when $c\ge0$.
\end{example}

\subsection{Fredholm alternative}

\begin{definition}[Formal adjoint operator and adjoint problem]
\label{def:formal-adjoint-operator}
\uses{def:divergence-nondivergence-form, def:bilinear-form-weak-solution-l2}
\dcref{ch6:6.2-def-1}
\group{Weak Solutions}

(i) The operator $L^*$, the formal adjoint of $L$, is
$$L^*v := -\sum_{i,j=1}^n(a^{ij}v_{x_j})_{x_i} - \sum_{i=1}^n b^iv_{x_i} + \left(c - \sum_{i=1}^n b^i_{x_i}\right)v,$$
provided $b^i \in C^1(\bar U)$ $(i=1,\ldots,n)$.

(ii) The adjoint bilinear form
$$B^* : H \times H \to \R$$
is defined by
$$B^*[v,u] = B[u,v]$$
for all $u,v \in H_0^1(U)$.

(iii) We say that $v \in H_0^1(U)$ is a weak solution of the adjoint problem
$$\begin{cases} L^*v = f & \text{in } U \\ v = 0 & \text{on } \partial U, \end{cases}$$
provided
$$B^*[v,u] = (f,u)$$
for all $u \in H_0^1(U)$.
\end{definition}

\begin{theorem}[Second Existence Theorem for weak solutions]
\label{thm:second-existence-fredholm-alternative}
\uses{thm:first-existence-weak-solutions, def:formal-adjoint-operator, thm:rellich-kondrachov}
\dcref{ch6:6.2-thm-4}
\group{Weak Solutions}

(i) Precisely one of the following statements holds:

either
$$\begin{cases} \text{for each } f \in L^2(U) \text{ there exists a unique} \\ \text{weak solution } u \text{ of the boundary-value problem} \\ (10) \quad \begin{cases} Lu = f & \text{in } U \\ u = 0 & \text{on } \partial U \end{cases} \end{cases} (\alpha)$$

or else
$$\begin{cases} \text{there exists a weak solution } u \not\equiv 0 \text{ of} \\ \text{the homogeneous problem} \\ (11) \quad \begin{cases} Lu = 0 & \text{in } U \\ u = 0 & \text{on } \partial U \end{cases} \end{cases} (\beta)$$

(ii) Furthermore, should assertion $(\beta)$ hold, the dimension of the subspace $N \subset H_0^1(U)$ of weak solutions of (11) is finite and equals the dimension of the subspace $N^* \subset H_0^1(U)$ of weak solutions of
$$(12) \qquad \begin{cases} L^*v = 0 & \text{in } U \\ v = 0 & \text{on } \partial U \end{cases}$$

(iii) Finally, the boundary-value problem (10) has a weak solution if and only if
$$(f,v) = 0 \text{ for all } v \in N^*.$$

The dichotomy $(\alpha)$, $(\beta)$ is the \textit{Fredholm alternative}.
\end{theorem}

\begin{proof}
\sketch Let $L_\gamma=L+\gamma I$ be the coercive isomorphism and set $K=\gamma L_\gamma^{-1}$ on $L^2(U)$. Rellich compactness makes $K$ compact, while $Lu=f$ is equivalent to $(I-K)u=L_\gamma^{-1}f$. The compact Fredholm alternative identifies the kernels of $I-K$ and its adjoint and gives the stated orthogonality condition.
\end{proof}

\begin{theorem}[Third Existence Theorem for weak solutions]
\label{thm:third-existence-spectrum}
\uses{thm:second-existence-fredholm-alternative, thm:energy-estimates}
\dcref{ch6:6.2-thm-5}
\group{Weak Solutions}

(i) There exists an at most countable set $\Sigma \subset \R$ such that the boundary-value problem
$$(24) \qquad \begin{cases} Lu = \lambda u + f & \text{in } U \\ u = 0 & \text{on } \partial U \end{cases}$$
has a unique weak solution for each $f \in L^2(U)$ if and only if $\lambda \notin \Sigma$.

(ii) If $\Sigma$ is infinite, then $\Sigma = \{\lambda_k\}_{k=1}^\infty$, the values of a nondecreasing sequence with
$$\lambda_k \to +\infty.$$
\end{theorem}

\begin{proof}
\sketch With $K=\gamma L_\gamma^{-1}$, the homogeneous equation $Lu=\lambda u$ is equivalent to $Ku=\gamma(\gamma+\lambda)^{-1}u$. Nonzero eigenvalues of compact $K$ form an at most countable set accumulating only at zero, which translates into an at most countable real spectrum tending to $+\infty$ when infinite.
\end{proof}

\begin{definition}[Spectrum of the operator $L$]
\label{def:spectrum-of-operator}
\uses{thm:third-existence-spectrum}
\dcref{ch6:6.2-def-2}
\group{Weak Solutions}

We call $\Sigma$ the (real) spectrum of the operator $L$.
\end{definition}

\begin{remark}[Eigenvalues and Helmholtz's equation]
\label{rem:eigenvalues-helmholtz-equation}
\uses{def:spectrum-of-operator}
\dcref{ch6:6.2-rem-3}
\group{Weak Solutions}

The homogeneous problem
$$\begin{cases} Lu = \lambda u & \text{in } U \\ u = 0 & \text{on } \partial U \end{cases}$$
has a nontrivial solution iff $\lambda\in\Sigma$; such a $\lambda$ is an eigenvalue and $w$ an eigenfunction. For $L=-\Delta$ this is the Helmholtz equation.
\end{remark}

\begin{theorem}[Boundedness of the inverse]
\label{thm:boundedness-of-inverse}
\uses{thm:third-existence-spectrum, thm:energy-estimates}
\dcref{ch6:6.2-thm-6}
\group{Weak Solutions}

If $\lambda \notin \Sigma$, there exists a constant $C$ such that
$$(29) \qquad \|u\|_{L^2(U)} \leq C\|f\|_{L^2(U)},$$
whenever $f \in L^2(U)$ and $u \in H_0^1(U)$ is the unique weak solution of
$$\begin{cases} Lu = \lambda u + f & \text{in } U \\ u = 0 & \text{on } \partial U. \end{cases}$$
The constant $C$ depends only on $\lambda$, $U$ and the coefficients of $L$.

No uniform choice of $C$ is possible as $\lambda$ approaches an eigenvalue.
\end{theorem}

\begin{proof}
\sketch If the estimate failed, normalize solutions with $\|u_k\|_{L^2}=1$ and right-hand sides tending to zero. Energy bounds and Rellich compactness yield a subsequence converging strongly in $L^2$ to a normalized homogeneous solution, contradicting $\lambda\notin\Sigma$.
\end{proof}

\begin{remark}[Complex solutions]
\label{rem:complex-valued-weak-solutions}
\uses{thm:lax-milgram-theorem, thm:second-existence-fredholm-alternative}
\dcref{ch6:6.2-rem-4}
\group{Weak Solutions}

For complex-valued $u,v\in H^1(U)$, set
$$(u,v)_{L^2(U)} := \int_U u\bar v\,dx, \quad (u,v)_{H^1(U)} := \int_U Du \cdot D\bar v + u\bar v\,dx,$$
and set
$$B[u,v] := \int_U \sum_{i,j=1}^n a^{ij}u_{x_i}\bar v_{x_j} + \sum_{i=1}^n b^iu_{x_i}\bar v + cu\bar v\,dx,$$
where $\bar{\ }$ denotes conjugation. Then
$$|B[u,v]| \leq \alpha\|u\|_{H_0^1(U)}\|v\|_{H_0^1(U)},$$
$$\beta\|u\|_{H_0^1(U)}^2 \leq \operatorname{Re}B[u,u] + \gamma\|u\|_{L^2(U)}^2 \quad (u,v \in H_0^1(U))$$
for suitable $\alpha,\beta>0$ and $\gamma\ge0$, and the complex Lax--Milgram/Fredholm theory gives the corresponding existence results.
\end{remark}

\section{Regularity}
\label{sec:elliptic-regularity}

\subsection{Interior regularity}

\begin{theorem}[Interior $H^2$-regularity]
\label{thm:interior-h2-regularity}
\uses{def:uniform-ellipticity, def:weak-solution-elliptic-bvp, thm:difference-quotients-weak-derivatives}
\dcref{ch6:6.3-thm-1}
\group{Regularity Theory}

Assume
$$(5) \qquad a^{ij} \in C^1(U), \quad b^i,c \in L^\infty(U) \quad (i,j=1,\ldots,n)$$
and
$$(6) \qquad f \in L^2(U).$$
Suppose furthermore that $u \in H^1(U)$ is a weak solution of the elliptic PDE
$$Lu = f \quad \text{in } U.$$
Then
$$(7) \qquad u \in H^2_{\text{loc}}(U);$$
and for each open $V \subset\subset U$ we have the estimate
$$(8) \qquad \|u\|_{H^2(V)} \leq C\left(\|f\|_{L^2(U)} + \|u\|_{L^2(U)}\right),$$
the constant $C$ depending only on $V$, $U$, and the coefficients of $L$.
\end{theorem}

\begin{proof}
\sketch Choose nested interior domains and test the weak equation with $-D_k^{-h}(\zeta^2D_k^hu)$. Ellipticity and Cauchy inequalities give a uniform $L^2$ bound for $D_k^hDu$; the difference-quotient criterion yields $H^2_{\mathrm{loc}}$. A second cutoff energy estimate replaces the $H^1$ term by $\|f\|_{L^2}+\|u\|_{L^2}$.
\end{proof}

\begin{remark}
\label{rem:interior-regularity-remarks}
\uses{thm:interior-h2-regularity}
\dcref{ch6:6.3-rem-1}
\group{Regularity Theory}

(i) No zero-trace condition is required: $u\in H^1(U)$ suffices.

(ii) Since $u\in H^2_{\mathrm{loc}}(U)$,
$$Lu = f \quad \text{a.e. in } U.$$
the weak identity agrees with the distributional integration-by-parts identity.
\end{remark}

\begin{theorem}[Higher interior regularity]
\label{thm:higher-interior-regularity}
\uses{thm:interior-h2-regularity}
\dcref{ch6:6.3-thm-2}
\group{Regularity Theory}

Let $m$ be a nonnegative integer, and assume
$$(25) \qquad a^{ij}, b^i, c \in C^{m+1}(U) \quad (i,j=1,\ldots,n)$$
and
$$(26) \qquad f \in H^m(U).$$
Suppose $u \in H^1(U)$ is a weak solution of the elliptic PDE
$$Lu = f \quad \text{in } U.$$
Then
$$(27) \qquad u \in H^{m+2}_{\text{loc}}(U);$$
and for each $V \subset\subset U$ we have the estimate
$$(28) \qquad \|u\|_{H^{m+2}(V)} \leq C(\|f\|_{H^m(U)} + \|u\|_{L^2(U)}),$$
the constant $C$ depending only on $m, U, V$ and the coefficients of $L$.
\end{theorem}

\begin{proof}
\sketch Induct on $m$. Differentiate the equation weakly, regard each first derivative of $u$ as solving an elliptic equation whose right-hand side is controlled by the induction hypothesis, apply interior $H^2$ regularity on nested domains, and sum the resulting estimates.
\end{proof}

\begin{theorem}[Infinite differentiability in the interior]
\label{thm:interior-infinite-differentiability}
\uses{thm:higher-interior-regularity, thm:general-sobolev-inequalities}
\dcref{ch6:6.3-thm-3}
\group{Regularity Theory}

Assume
$$a^{ij}, b^i, c \in C^\infty(U) \quad (i,j=1,\ldots,n)$$
and
$$f \in C^\infty(U).$$
Suppose $u \in H^1(U)$ is a weak solution of the elliptic PDE
$$Lu = f \quad \text{in } U.$$
Then
$$u \in C^\infty(U).$$
No boundary condition on $u$ is required; boundary singularities do not affect the interior regularity.
\end{theorem}

\begin{proof}
\sketch Higher interior regularity places $u$ in every local Sobolev space. Sobolev embedding then gives derivatives of every order on compact subsets of $U$.
\end{proof}

\subsection{Boundary regularity}

\begin{theorem}[Boundary $H^2$-regularity]
\label{thm:boundary-h2-regularity}
\uses{def:uniform-ellipticity, def:weak-solution-elliptic-bvp, thm:difference-quotients-weak-derivatives, thm:interior-h2-regularity}
\dcref{ch6:6.3-thm-4}
\group{Regularity Theory}

Assume
$$(38) \qquad a^{ij} \in C^1(\bar U), \quad b^i,c \in L^\infty(U) \quad (i,j=1,\ldots,n)$$
and
$$(39) \qquad f \in L^2(U).$$
Suppose that $u \in H_0^1(U)$ is a weak solution of the elliptic boundary-value problem
$$(40) \qquad \begin{cases} Lu = f & \text{in } U \\ u = 0 & \text{on } \partial U. \end{cases}$$
Assume finally
$$(41) \qquad \partial U \text{ is } C^2.$$
Then
$$u \in H^2(U),$$
and we have the estimate
$$(42) \qquad \|u\|_{H^2(U)} \leq C(\|f\|_{L^2(U)}+\|u\|_{L^2(U)}),$$
the constant $C$ depending only on $U$ and the coefficients of $L$.
\end{theorem}

\begin{proof}
\sketch Flatten the boundary and apply tangential difference-quotient tests to obtain all tangential second derivatives with uniform estimates. The equation and uniform ellipticity solve for the remaining double normal derivative; interior regularity and a finite boundary cover give the global estimate.
\end{proof}

\begin{remark}
\label{rem:boundary-regularity-remarks}
\uses{thm:boundary-h2-regularity, thm:boundedness-of-inverse}
\dcref{ch6:6.3-rem-2}
\group{Regularity Theory}

(i) For the unique zero-trace solution of (40),
$$\|u\|_{H^2(U)} \leq C\|f\|_{L^2(U)}.$$

(ii) The estimate uses the trace condition $u=0$ on $\partial U$.
\end{remark}

\begin{theorem}[Higher boundary regularity]
\label{thm:higher-boundary-regularity}
\uses{thm:boundary-h2-regularity, thm:higher-interior-regularity}
\dcref{ch6:6.3-thm-5}
\group{Regularity Theory}

Let $m$ be a nonnegative integer, and assume
$$(72) \qquad a^{ij}, b^i, c \in C^{m+1}(\bar U) \quad (i,j=1,\ldots,n)$$
and
$$(73) \qquad f \in H^m(U).$$
Suppose that $u \in H_0^1(U)$ is a weak solution of the boundary-value problem
$$(74) \qquad \begin{cases} Lu = f & \text{in } U \\ u = 0 & \text{on } \partial U \end{cases}$$
Assume finally
$$(75) \qquad \partial U \text{ is } C^{m+2}.$$
Then
$$(76) \qquad u \in H^{m+2}(U),$$
and we have the estimate
$$(77) \qquad \|u\|_{H^{m+2}(U)} \leq C(\|f\|_{H^m(U)} + \|u\|_{L^2(U)}),$$
the constant $C$ depending only on $m, U$ and the coefficients of $L$.
\end{theorem}

\begin{proof}
\sketch Induct on $m$ after flattening each boundary chart. Tangentially differentiate the transformed equation, apply the lower-order boundary estimate to the differentiated equations, and recover derivatives with additional normal components from the elliptic equation. Patch the chart estimates with interior regularity.
\end{proof}

\begin{remark}
\label{rem:higher-boundary-regularity-simplified}
\uses{thm:higher-boundary-regularity}
\dcref{ch6:6.3-rem-3}
\group{Regularity Theory}

For the unique solution of (74),
$$\|u\|_{H^{m+2}(U)} \leq C\|f\|_{H^m(U)}.$$
\end{remark}

\begin{theorem}[Infinite differentiability up to the boundary]
\label{thm:boundary-infinite-differentiability}
\uses{thm:higher-boundary-regularity, thm:general-sobolev-inequalities}
\dcref{ch6:6.3-thm-6}
\group{Regularity Theory}

Assume
$$a^{ij}, b^i, c \in C^\infty(\bar U) \quad (i,j=1,\ldots,n)$$
and
$$f \in C^\infty(\bar U).$$
Suppose $u \in H_0^1(U)$ is a weak solution of the boundary-value problem
$$\begin{cases} Lu = f & \text{in } U \\ u = 0 & \text{on } \partial U. \end{cases}$$
Assume also that $\partial U$ is $C^\infty$. Then
$$u \in C^\infty(\bar U).$$
\end{theorem}

\begin{proof}
\sketch Higher boundary regularity yields $H^k(U)$ for every $k$. Sobolev embedding on the smooth bounded domain gives $C^\infty(\bar U)$.
\end{proof}

\section{Maximum Principles}
\label{sec:elliptic-maximum-principles}

\subsection{Weak maximum principle}

\begin{theorem}[Weak maximum principle]
\label{thm:weak-maximum-principle-c-zero}
\uses{def:uniform-ellipticity, def:divergence-nondivergence-form, thm:boundary-h2-regularity}
\dcref{ch6:6.4-thm-1}
\group{Maximum Principles}

Assume $u \in C^2(U) \cap C(\bar U)$ and
$$c \equiv 0 \quad \text{in } U.$$

(i) If
$$Lu \leq 0 \quad \text{in } U,$$
then
$$\max_{\bar U} u = \max_{\partial U} u.$$

(ii) If
$$Lu \geq 0 \quad \text{in } U,$$
then
$$\min_{\bar U} u = \min_{\partial U} u.$$
\end{theorem}

\begin{proof}
\sketch At a strict interior maximum, $Du=0$ and $D^2u$ is negative semidefinite; diagonalizing the positive coefficient matrix gives $Lu\ge0$, contradicting $Lu<0$. For the nonstrict case perturb by $\varepsilon e^{\lambda x_1}$ with $\lambda$ large and let $\varepsilon\downarrow0$; apply the result to $-u$ for minima.
\end{proof}

\begin{remark}
\label{rem:subsolution-supersolution-terminology}
\uses{thm:weak-maximum-principle-c-zero}
\dcref{ch6:6.4-rem-1}
\group{Maximum Principles}

A function satisfying (3) is a subsolution and attains its maximum on $\partial U$; (4) defines a supersolution attaining its minimum there.
\end{remark}

\begin{theorem}[Weak maximum principle for $c \geq 0$]
\label{thm:weak-maximum-principle-c-nonneg}
\uses{thm:weak-maximum-principle-c-zero}
\dcref{ch6:6.4-thm-2}
\group{Maximum Principles}

Assume $u \in C^2(U) \cap C(\bar U)$ and
$$c \geq 0 \quad \text{in } U.$$

(i) If
$$Lu \leq 0 \quad \text{in } U,$$
then
$$(11) \qquad \max_{\bar U} u \leq \max_{\partial U} u^+.$$

(ii) Likewise, if
$$Lu \geq 0 \quad \text{in } U,$$
then
$$(12) \qquad \min_{\bar U} u \geq -\max_{\partial U} u^-.$$
\end{theorem}

\begin{proof}
\sketch On $V=\{u>0\}$, the operator $K=L-c$ has no zeroth-order term and satisfies $Ku\le0$. Apply the zero-order-free maximum principle on $V$ and then apply the result to $-u$.
\end{proof}

\begin{remark}
\label{rem:weak-maximum-principle-absolute-value}
\uses{thm:weak-maximum-principle-c-nonneg}
\dcref{ch6:6.4-rem-2}
\group{Maximum Principles}

If $Lu=0$ in $U$, then
$$(13) \qquad \max_{\bar U}|u| = \max_{\partial U}|u|.$$
\end{remark}

\subsection{Strong maximum principle}

\begin{lemma}[Hopf's Lemma]
\label{lem:hopfs-lemma}
\uses{def:uniform-ellipticity, thm:weak-maximum-principle-c-zero}
\dcref{ch6:6.4-lem-1}
\group{Maximum Principles}

Assume $u \in C^2(U) \cap C^1(\bar U)$ and
$$c = 0 \quad \text{in } U.$$
Suppose further
$$Lu \leq 0 \quad \text{in } U,$$
and there exists a point $x^0 \in \partial U$ such that
$$(14) \qquad u(x^0) > u(x) \quad \text{for all } x \in U.$$
Assume finally that $U$ satisfies the interior ball condition at $x^0$; that is, there exists an open ball $B \subset U$ with $x^0 \in \partial B$.

(i) Then
$$\frac{\partial u}{\partial\nu}(x^0) > 0,$$
where $\nu$ is the outer unit normal to $B$ at $x^0$.

(ii) If
$$c \geq 0 \quad \text{in } U,$$
the same conclusion holds provided
$$u(x^0) \geq 0.$$
\end{lemma}

\begin{proof}
\sketch Use the exponential radial barrier $v=e^{-\lambda|x|^2}-e^{-\lambda r^2}$ on an interior annulus. For large $\lambda$, ellipticity gives $Lv\le0$; comparison with $u+\varepsilon v-u(x^0)$ and differentiation at $x^0$ yield a strictly positive outer normal derivative.
\end{proof}

\begin{remark}
\label{rem:hopfs-lemma-strict-inequality}
\uses{lem:hopfs-lemma}
\dcref{ch6:6.4-rem-3}
\group{Maximum Principles}

Part (i) strengthens the obvious non-strict inequality to $\frac{\partial u}{\partial\nu}(x^0)>0$. The interior ball condition holds when $\partial U$ is $C^2$.
\end{remark}

\begin{theorem}[Strong maximum principle]
\label{thm:strong-maximum-principle-c-zero}
\uses{lem:hopfs-lemma}
\dcref{ch6:6.4-thm-3}
\group{Maximum Principles}

Assume $u \in C^2(U) \cap C(\bar U)$ and
$$c \equiv 0 \quad \text{in } U.$$
Suppose also $U$ is connected, open and bounded.

(i) If
$$Lu \leq 0 \quad \text{in } U$$
and $u$ attains its maximum over $\bar U$ at an interior point, then
$$u \text{ is constant within } U.$$

(ii) Similarly, if
$$Lu \geq 0 \quad \text{in } U$$
and $u$ attains its minimum over $\bar U$ at an interior point, then
$$u \text{ is constant within } U.$$
\end{theorem}

\begin{proof}
\sketch If a nonconstant subsolution attained its maximum inside, take a largest ball in the strict sublevel set tangent to the maximum set. Hopf lemma gives a nonzero normal derivative at the tangency point, contradicting vanishing of the gradient at an interior maximum.
\end{proof}

\begin{theorem}[Strong maximum principle with $c \geq 0$]
\label{thm:strong-maximum-principle-c-nonneg}
\uses{lem:hopfs-lemma, thm:strong-maximum-principle-c-zero}
\dcref{ch6:6.4-thm-4}
\group{Maximum Principles}

Assume $u \in C^2(U) \cap C(\bar U)$ and
$$c \geq 0 \quad \text{in } U.$$
Suppose also $U$ is connected.

(i) If
$$Lu \leq 0 \quad \text{in } U$$
and $u$ attains a nonnegative maximum over $\bar U$ at an interior point, then
$$u \text{ is constant within } U.$$

(ii) Similarly, if
$$Lu \geq 0 \quad \text{in } U$$
and $u$ attains a nonpositive minimum over $\bar U$ at an interior point, then
$$u \text{ is constant within } U.$$

\end{theorem}

\begin{proof}
\sketch Repeat the tangency-ball argument for a nonnegative maximum and use part (ii) of Hopf lemma; apply the result to $-u$ for minima.
\end{proof}

\subsection{Harnack's inequality}

\begin{theorem}[Harnack's inequality]
\label{thm:harnacks-inequality-elliptic}
\dcref{ch6:6.4-thm-5}
\group{Maximum Principles}
Assume $u \geq 0$ is a $C^2$ solution of
$$Lu = 0 \quad \text{in } U,$$
and suppose $V \subset\subset U$ is connected. Then there exists a constant $C$ such that
$$(18) \qquad \sup_V u \leq C\inf_V u.$$
The constant $C$ depends only on $V$ and the coefficients of $L$.
\end{theorem}
For bounded measurable coefficient variants, see \cite{GilbargTrudinger1983}.

\section{Eigenvalues and Eigenfunctions}
\label{sec:elliptic-eigenvalues}
Perron--Frobenius comparison: \cite{Gantmacher1977}.

\subsection{The symmetric case}

\begin{theorem}[Eigenvalues of symmetric elliptic operators]
\label{thm:existence-eigenvalues-symmetric-operator}
\uses{thm:second-existence-fredholm-alternative, thm:rellich-kondrachov, def:bilinear-form-weak-solution-l2}
\dcref{ch6:6.5-thm-1}
\group{Eigenvalues}

(i) Each eigenvalue of $L$ is real.

(ii) Repeating each eigenvalue according to its finite multiplicity gives
$$\Sigma = \{\lambda_k\}_{k=1}^\infty,$$
where
$$0 < \lambda_1 \leq \lambda_2 \leq \lambda_3 \leq \cdots,$$
and
$$\lambda_k \to \infty \quad \text{as } k \to \infty.$$

(iii) Finally, there exists an orthonormal basis $\{w_k\}_{k=1}^\infty$ of $L^2(U)$, where $w_k \in H_0^1(U)$ is an eigenfunction corresponding to $\lambda_k$:
$$(4) \qquad \begin{cases} Lw_k = \lambda_k w_k & \text{in } U \\ w_k = 0 & \text{on } \partial U \end{cases}$$
for $k = 1,2,\ldots$.
\end{theorem}

\begin{proof}
\sketch The inverse $S=L^{-1}$ is compact on $L^2(U)$ and symmetry of $B$ makes $S$ self-adjoint and positive. The compact self-adjoint spectral theorem gives a real positive eigenbasis for $S$; reciprocating its eigenvalues gives the asserted sequence for $L$.
\end{proof}

\begin{remark}
\label{rem:eigenfunctions-smoothness}
\uses{thm:existence-eigenvalues-symmetric-operator, thm:interior-infinite-differentiability, thm:boundary-infinite-differentiability}
\dcref{ch6:6.5-rem-1}
\group{Eigenvalues}

Each $w_k\in C^\infty(U)$, and $w_k\in C^\infty(\bar U)$ when $\partial U$ is smooth.
\end{remark}

\begin{definition}[Principal eigenvalue]
\label{def:principal-eigenvalue}
\uses{thm:existence-eigenvalues-symmetric-operator}
\dcref{ch6:6.5-def-1}
\group{Eigenvalues}

We call $\lambda_1 > 0$ the principal eigenvalue of $L$.
\end{definition}

\begin{theorem}[Variational principle for the principal eigenvalue]
\label{thm:variational-principle-principal-eigenvalue}
\uses{thm:existence-eigenvalues-symmetric-operator, def:principal-eigenvalue, thm:strong-maximum-principle-c-zero}
\dcref{ch6:6.5-thm-2}
\group{Eigenvalues}

(i) We have
$$(5) \qquad \lambda_1 = \min\{B[u,u] \mid u \in H_0^1(U), \|u\|_{L^2} = 1\}.$$

(ii) The minimum is attained by a function $w_1>0$ in $U$ solving
$$\begin{cases} Lw_1 = \lambda_1w_1 & \text{in } U \\ w_1 = 0 & \text{on } \partial U. \end{cases}$$

(iii) Finally, if $u \in H_0^1(U)$ is any weak solution of
$$\begin{cases} Lu = \lambda_1u & \text{in } U \\ u = 0 & \text{on } \partial U, \end{cases}$$
then $u$ is a multiple of $w_1$.
\end{theorem}

\begin{proof}
\sketch Expand normalized $u$ in the orthonormal eigenbasis to obtain $B[u,u]=\sum d_k^2\lambda_k\ge\lambda_1$. Equality forces membership in the first eigenspace. Applying this to $u^+$ and $u^-$ and then the strong maximum principle gives a one-signed eigenfunction; a zero-integral linear combination proves simplicity.
\end{proof}

\begin{remark}[Simplicity of $\lambda_1$ and Rayleigh's formula]
\label{rem:principal-eigenvalue-simple-rayleigh}
\uses{thm:variational-principle-principal-eigenvalue}
\dcref{ch6:6.5-rem-2}
\group{Eigenvalues}

(i) The principal eigenvalue is simple, and
$$0 < \lambda_1 < \lambda_2 \leq \lambda_3 \leq \cdots.$$

(ii) The equivalent Rayleigh formula is
$$\lambda_1 = \min_{\substack{u \in H_0^1(U) \\ u \neq 0}} \frac{B[u,u]}{\|u\|_{L^2(U)}^2}.$$
\end{remark}

\subsection{Eigenvalues of nonsymmetric elliptic operators}

\begin{theorem}[Principal eigenvalue for nonsymmetric elliptic operators]
\label{thm:principal-eigenvalue-nonsymmetric-operator}
\uses{thm:strong-maximum-principle-c-nonneg, lem:hopfs-lemma, thm:weak-maximum-principle-c-nonneg}
\dcref{ch6:6.5-thm-3}
\group{Eigenvalues}

(i) There exists a real eigenvalue $\lambda_1$ for the operator $L$, taken with zero boundary conditions. Furthermore if $\lambda \in \C$ is any other eigenvalue, we have
$$\operatorname{Re}(\lambda) \geq \lambda_1.$$

(ii) There exists a corresponding eigenfunction $w_1$, which is positive within $U$.

(iii) The eigenvalue $\lambda_1$ is simple; that is, if $u$ is any solution of (1), then $u$ is a multiple of $w_1$.
\end{theorem}

\begin{proof}
\sketch Apply the compact positive inverse to the cone of nonnegative functions; the Schaefer continuation argument produces a positive eigenfunction $w_1$. For any complex eigenfunction, apply the maximum principle to $|u/w_1^{1-\varepsilon}|^2$ and let $\varepsilon\downarrow0$ to obtain $\operatorname{Re}\lambda\ge\lambda_1$. Strong maximum principle and Hopf lemma applied to $w_1-\chi u$ give simplicity.
\end{proof}

\section{References}
\label{sec:elliptic-references}
\textbf{Elliptic foundations.} \cite{GilbargTrudinger1983}.
\textbf{Regularity.} \cite{GilbargTrudinger1983}, \cite{Krylov1996}, \cite{LadyzhenskayaUraltseva1968}.
\textbf{Maximum principles.} \cite{GilbargTrudinger1983}, \cite{ProtterWeinberger1967}.
\textbf{Eigenvalues.} \cite{Smoller1983}, \cite{ProtterWeinberger1967}.
